% assets/w09-arc-bank.svg — 호 다발 21개 중 하나를 고르는 법 (collision_avoidance_enu.py)
%
%   bash _tools/tikz2svg.sh assets/src/w09-arc-bank.tex assets/w09-arc-bank.svg
%
% 이 그림이 말하는 것 하나
%   회피기는 "어디로 갈까" 를 자유롭게 고르지 않는다. **배가 실제로 낼 수 있는 선회율**
%   21개를 미리 깔아 두고 (-12 ~ +12 deg/s), 각각을 3초 앞까지 따라가 보고, 부딪히는 것을
%   지우고 남은 것 중 점수가 가장 높은 하나를 고른다. 3초는 4.5 m 다 — 자기 배 길이(4.9 m)
%   보다 짧다. 그래서 이 회피기는 멀리 보는 기계가 아니라 가까이를 놓치지 않는 기계다.
%
% 화면 : 위가 선수(x_b), 오른쪽이 우현. 즉 배를 위에서 본 그림이다. 1 m = 20 mm
%   배치 규칙 — 그림(호·선체·장애물)은 \S 로 재고, 글상자는 mm 절대좌표로 둔다.
%   글상자까지 \S 로 재면 축척을 바꿀 때마다 글이 그림 밖으로 날아간다 (2026-10-27 에 겪었다)
%
% 호의 식 : 선회율 r, 속력 u 로 돌면 반지름 R = u/r, 시간 T 뒤 선회각 Phi = r*T
%   u=1.5, r=12 deg/s -> R = 7.16 m, Phi = 36 deg, 호 길이 4.5 m, 끝점 (-1.368, 4.209)
%   좌선회는 화면 왼쪽으로 휘므로 가로 성분에 음부호가 붙는다 (2026-10-27 에 한 번 틀렸다)
%
% 함정 둘 — 둘 다 2026-10-27 에 걸렸다
%   (1) 줄바꿈은 두 겹 역슬래시만 쓴다. 길이를 달면(`[2pt]`) kotex/standalone 조합에서
%       수식 시작으로 읽혀 "Command \scriptsize invalid in math mode" 가 난다
%   (2) 이 파일은 **Write 도구로** 쓴다. 셸 히어독을 거치면 두 겹 역슬래시가 한 겹으로
%       줄어 줄바꿈이 통째로 사라진다 — 오류 없이 한 줄로 붙어 나오므로 더 나쁘다

\documentclass[border=6pt]{standalone}
% capstone-style.tex — 이 볼트의 TikZ 그림이 공유하는 스타일
%
% 쓰는 법:  % capstone-style.tex — 이 볼트의 TikZ 그림이 공유하는 스타일
%
% 쓰는 법:  % capstone-style.tex — 이 볼트의 TikZ 그림이 공유하는 스타일
%
% 쓰는 법:  \input{capstone-style}   (assets/src/ 안에서)
% 빌드:     bash _tools/tikz2svg.sh assets/src/w02-node-topic.tex assets/w02-node-topic.svg
%
% 왜 TikZ 인가
%   손으로 SVG 좌표를 쓰면 화살촉·선 굵기·글자 크기가 그림마다 미묘하게 달라진다.
%   그것이 "자료가 어설퍼 보인다"의 정체다. 스타일을 한 번만 정하고 상대좌표로
%   배치하면 좌표를 고쳐도 그림이 무너지지 않는다.
%   수식은 본문과 같은 TeX 글꼴로 나온다.
%
% 한글
%   ko.TeX(DVI 경로)를 쓴다. XeLaTeX 은 TikZ 그래픽을 PDF special 로 내보내는데
%   dvisvgm 이 그것을 받으려면 Ghostscript < 10.01 이 필요하다. 이 PC 에는
%   10.07 이 깔려 있어 그 경로가 막힌다. latex -> DVI -> dvisvgm 이 유일하게
%   동작하는 조합이다.

\usepackage[hangul]{kotex}
\usepackage{tikz}
\usepackage{amsmath}
\usepackage{xcolor}
\usetikzlibrary{arrows.meta, positioning, calc, fit, backgrounds, decorations.pathmorphing}

% ---- 색 --------------------------------------------------------------------
% 강조는 하나만. 나머지는 검정. 흑백 인쇄에서도 읽혀야 한다.
\definecolor{ink}{HTML}{1A1A1A}
\definecolor{soft}{HTML}{666666}
\definecolor{accent}{HTML}{7C3AED}
\definecolor{warn}{HTML}{B91C1C}

% 계층 색 — Simulink 모델의 블록 색과 같은 규칙을 쓴다
\definecolor{guid}{HTML}{1D4ED8}    % 유도
\definecolor{ctrl}{HTML}{EA580C}    % 제어
\definecolor{thr}{HTML}{D97706}     % 추진기
\definecolor{plant}{HTML}{16A34A}   % 운동모델
\definecolor{nav}{HTML}{7C3AED}     % 항법 · 신호처리
\definecolor{pathc}{HTML}{0D9488}   % 계획 경로 · 항적
\definecolor{errc}{HTML}{C2410C}    % 오차
\definecolor{flow}{HTML}{0369A1}    % 조류 · 바람

% ---- 선과 화살촉 -----------------------------------------------------------
\tikzset{
  sig/.style   = {ink, line width=0.5pt, -{Stealth[length=4pt, width=3pt]}},
  wire/.style  = {ink, line width=0.5pt},
  fb/.style    = {ink, line width=0.5pt, densely dashed,
                  -{Stealth[length=4pt, width=3pt]}},
  asig/.style  = {accent, line width=0.5pt, densely dashed,
                  -{Stealth[length=4pt, width=3pt]}},
  % 블록 — 흰 바탕, 직각, 검은 테두리. 색은 테두리에만 준다
  blk/.style   = {draw=ink, line width=0.55pt, fill=white,
                  minimum width=18mm, minimum height=8.5mm, inner sep=2pt, font=\small},
  gblk/.style  = {blk, draw=guid},
  cblk/.style  = {blk, draw=ctrl},
  tblk/.style  = {blk, draw=thr},
  pblk/.style  = {blk, draw=plant},
  nblk/.style  = {blk, draw=nav},
  % 합산점 · 분기점
  sum/.style   = {draw=ink, line width=0.55pt, fill=white, circle,
                  inner sep=0pt, minimum size=4.6mm, font=\scriptsize},
  dot/.style   = {ink, fill=ink, circle, inner sep=0pt, minimum size=1.5mm},
  % 글자
  sname/.style = {font=\small, inner sep=1.5pt},
  note/.style  = {font=\scriptsize, soft, inner sep=1.5pt},
  wnote/.style = {font=\scriptsize, warn, inner sep=1.5pt},
  anote/.style = {font=\scriptsize, accent, inner sep=1.5pt},
  axis/.style  = {ink, line width=0.5pt},
  tick/.style  = {font=\scriptsize, soft},
  % 묶음 상자 — 서브시스템 표시
  grp/.style   = {draw=soft, line width=0.4pt, densely dotted, rounded corners=1mm,
                  inner sep=3mm},
}
   (assets/src/ 안에서)
% 빌드:     bash _tools/tikz2svg.sh assets/src/w02-node-topic.tex assets/w02-node-topic.svg
%
% 왜 TikZ 인가
%   손으로 SVG 좌표를 쓰면 화살촉·선 굵기·글자 크기가 그림마다 미묘하게 달라진다.
%   그것이 "자료가 어설퍼 보인다"의 정체다. 스타일을 한 번만 정하고 상대좌표로
%   배치하면 좌표를 고쳐도 그림이 무너지지 않는다.
%   수식은 본문과 같은 TeX 글꼴로 나온다.
%
% 한글
%   ko.TeX(DVI 경로)를 쓴다. XeLaTeX 은 TikZ 그래픽을 PDF special 로 내보내는데
%   dvisvgm 이 그것을 받으려면 Ghostscript < 10.01 이 필요하다. 이 PC 에는
%   10.07 이 깔려 있어 그 경로가 막힌다. latex -> DVI -> dvisvgm 이 유일하게
%   동작하는 조합이다.

\usepackage[hangul]{kotex}
\usepackage{tikz}
\usepackage{amsmath}
\usepackage{xcolor}
\usetikzlibrary{arrows.meta, positioning, calc, fit, backgrounds, decorations.pathmorphing}

% ---- 색 --------------------------------------------------------------------
% 강조는 하나만. 나머지는 검정. 흑백 인쇄에서도 읽혀야 한다.
\definecolor{ink}{HTML}{1A1A1A}
\definecolor{soft}{HTML}{666666}
\definecolor{accent}{HTML}{7C3AED}
\definecolor{warn}{HTML}{B91C1C}

% 계층 색 — Simulink 모델의 블록 색과 같은 규칙을 쓴다
\definecolor{guid}{HTML}{1D4ED8}    % 유도
\definecolor{ctrl}{HTML}{EA580C}    % 제어
\definecolor{thr}{HTML}{D97706}     % 추진기
\definecolor{plant}{HTML}{16A34A}   % 운동모델
\definecolor{nav}{HTML}{7C3AED}     % 항법 · 신호처리
\definecolor{pathc}{HTML}{0D9488}   % 계획 경로 · 항적
\definecolor{errc}{HTML}{C2410C}    % 오차
\definecolor{flow}{HTML}{0369A1}    % 조류 · 바람

% ---- 선과 화살촉 -----------------------------------------------------------
\tikzset{
  sig/.style   = {ink, line width=0.5pt, -{Stealth[length=4pt, width=3pt]}},
  wire/.style  = {ink, line width=0.5pt},
  fb/.style    = {ink, line width=0.5pt, densely dashed,
                  -{Stealth[length=4pt, width=3pt]}},
  asig/.style  = {accent, line width=0.5pt, densely dashed,
                  -{Stealth[length=4pt, width=3pt]}},
  % 블록 — 흰 바탕, 직각, 검은 테두리. 색은 테두리에만 준다
  blk/.style   = {draw=ink, line width=0.55pt, fill=white,
                  minimum width=18mm, minimum height=8.5mm, inner sep=2pt, font=\small},
  gblk/.style  = {blk, draw=guid},
  cblk/.style  = {blk, draw=ctrl},
  tblk/.style  = {blk, draw=thr},
  pblk/.style  = {blk, draw=plant},
  nblk/.style  = {blk, draw=nav},
  % 합산점 · 분기점
  sum/.style   = {draw=ink, line width=0.55pt, fill=white, circle,
                  inner sep=0pt, minimum size=4.6mm, font=\scriptsize},
  dot/.style   = {ink, fill=ink, circle, inner sep=0pt, minimum size=1.5mm},
  % 글자
  sname/.style = {font=\small, inner sep=1.5pt},
  note/.style  = {font=\scriptsize, soft, inner sep=1.5pt},
  wnote/.style = {font=\scriptsize, warn, inner sep=1.5pt},
  anote/.style = {font=\scriptsize, accent, inner sep=1.5pt},
  axis/.style  = {ink, line width=0.5pt},
  tick/.style  = {font=\scriptsize, soft},
  % 묶음 상자 — 서브시스템 표시
  grp/.style   = {draw=soft, line width=0.4pt, densely dotted, rounded corners=1mm,
                  inner sep=3mm},
}
   (assets/src/ 안에서)
% 빌드:     bash _tools/tikz2svg.sh assets/src/w02-node-topic.tex assets/w02-node-topic.svg
%
% 왜 TikZ 인가
%   손으로 SVG 좌표를 쓰면 화살촉·선 굵기·글자 크기가 그림마다 미묘하게 달라진다.
%   그것이 "자료가 어설퍼 보인다"의 정체다. 스타일을 한 번만 정하고 상대좌표로
%   배치하면 좌표를 고쳐도 그림이 무너지지 않는다.
%   수식은 본문과 같은 TeX 글꼴로 나온다.
%
% 한글
%   ko.TeX(DVI 경로)를 쓴다. XeLaTeX 은 TikZ 그래픽을 PDF special 로 내보내는데
%   dvisvgm 이 그것을 받으려면 Ghostscript < 10.01 이 필요하다. 이 PC 에는
%   10.07 이 깔려 있어 그 경로가 막힌다. latex -> DVI -> dvisvgm 이 유일하게
%   동작하는 조합이다.

\usepackage[hangul]{kotex}
\usepackage{tikz}
\usepackage{amsmath}
\usepackage{xcolor}
\usetikzlibrary{arrows.meta, positioning, calc, fit, backgrounds, decorations.pathmorphing}

% ---- 색 --------------------------------------------------------------------
% 강조는 하나만. 나머지는 검정. 흑백 인쇄에서도 읽혀야 한다.
\definecolor{ink}{HTML}{1A1A1A}
\definecolor{soft}{HTML}{666666}
\definecolor{accent}{HTML}{7C3AED}
\definecolor{warn}{HTML}{B91C1C}

% 계층 색 — Simulink 모델의 블록 색과 같은 규칙을 쓴다
\definecolor{guid}{HTML}{1D4ED8}    % 유도
\definecolor{ctrl}{HTML}{EA580C}    % 제어
\definecolor{thr}{HTML}{D97706}     % 추진기
\definecolor{plant}{HTML}{16A34A}   % 운동모델
\definecolor{nav}{HTML}{7C3AED}     % 항법 · 신호처리
\definecolor{pathc}{HTML}{0D9488}   % 계획 경로 · 항적
\definecolor{errc}{HTML}{C2410C}    % 오차
\definecolor{flow}{HTML}{0369A1}    % 조류 · 바람

% ---- 선과 화살촉 -----------------------------------------------------------
\tikzset{
  sig/.style   = {ink, line width=0.5pt, -{Stealth[length=4pt, width=3pt]}},
  wire/.style  = {ink, line width=0.5pt},
  fb/.style    = {ink, line width=0.5pt, densely dashed,
                  -{Stealth[length=4pt, width=3pt]}},
  asig/.style  = {accent, line width=0.5pt, densely dashed,
                  -{Stealth[length=4pt, width=3pt]}},
  % 블록 — 흰 바탕, 직각, 검은 테두리. 색은 테두리에만 준다
  blk/.style   = {draw=ink, line width=0.55pt, fill=white,
                  minimum width=18mm, minimum height=8.5mm, inner sep=2pt, font=\small},
  gblk/.style  = {blk, draw=guid},
  cblk/.style  = {blk, draw=ctrl},
  tblk/.style  = {blk, draw=thr},
  pblk/.style  = {blk, draw=plant},
  nblk/.style  = {blk, draw=nav},
  % 합산점 · 분기점
  sum/.style   = {draw=ink, line width=0.55pt, fill=white, circle,
                  inner sep=0pt, minimum size=4.6mm, font=\scriptsize},
  dot/.style   = {ink, fill=ink, circle, inner sep=0pt, minimum size=1.5mm},
  % 글자
  sname/.style = {font=\small, inner sep=1.5pt},
  note/.style  = {font=\scriptsize, soft, inner sep=1.5pt},
  wnote/.style = {font=\scriptsize, warn, inner sep=1.5pt},
  anote/.style = {font=\scriptsize, accent, inner sep=1.5pt},
  axis/.style  = {ink, line width=0.5pt},
  tick/.style  = {font=\scriptsize, soft},
  % 묶음 상자 — 서브시스템 표시
  grp/.style   = {draw=soft, line width=0.4pt, densely dotted, rounded corners=1mm,
                  inner sep=3mm},
}


\begin{document}
\begin{tikzpicture}[x=1mm, y=1mm]

\def\S{20}        % 1 m = 20 mm

% 호 하나 — #1 반지름[m] (양수면 좌선회, 음수면 우선회), #2 선회각[deg], #3 스타일
\newcommand{\arcpath}[3]{%
  \draw[#3] plot[domain=0:#2, samples=28, variable=\t]
       ({-(#1)*\S*(1-cos(\t))}, {abs(#1)*\S*sin(\t)});}

% ===================== 호 다발 =====================
% 선회율 -12..+12 deg/s 를 21개로 나누면 간격 1.2 deg/s. 그림에는 둘 건너 하나씩 11개만
% 그린다 — 21개를 다 그리면 선이 붙어 읽히지 않는다 (코드는 21개 전부를 시험한다)

% (가) 우선회 다발 — 장애물 쪽. 3초 안에 선체가 닿는다
\arcpath{-7.16}{36}{warn, densely dotted, line width=0.6pt}     % r = -12.0 deg/s
\arcpath{-8.95}{28.8}{warn, densely dotted, line width=0.6pt}   % r =  -9.6
\arcpath{-11.94}{21.6}{warn, densely dotted, line width=0.6pt}  % r =  -7.2
\arcpath{-17.90}{14.4}{warn, densely dotted, line width=0.6pt}  % r =  -4.8
\arcpath{-35.81}{7.2}{warn, densely dotted, line width=0.6pt}   % r =  -2.4

% (나) 직진 — 목표로 바로 가는 길. 이것이 막혔는가가 oa_flag 의 질문이다
\draw[warn, densely dotted, line width=0.8pt] (0,0) -- (0,4.5*\S);

% (다) 좌선회 다발 — 열린 쪽. 살아남은 호
\arcpath{35.81}{7.2}{soft, line width=0.6pt}
\arcpath{17.90}{14.4}{soft, line width=0.6pt}
\arcpath{11.94}{21.6}{soft, line width=0.6pt}
\arcpath{8.95}{28.8}{soft, line width=0.6pt}

% (라) 고른 호 — 좌선회 최대 +12 deg/s
\arcpath{7.16}{36}{pathc, line width=1.8pt}
\fill[pathc] (-1.368*\S, 4.209*\S) circle (1.6);

% 3초 뒤의 배 — 36 deg 돌아 있다 (좌선회이므로 화면에서 반시계)
\begin{scope}[shift={(-1.368*\S, 4.209*\S)}, rotate=36]
  \draw[pathc, line width=0.7pt, densely dashed]
       (-1.23*\S, -2.4*\S) rectangle (1.23*\S, 2.5*\S);
\end{scope}
\node[note, pathc, anchor=east, text width=38mm, align=right] at (-90, 126)
     {3초 뒤의 배 \\ $36^\circ$ 돌았고 $4.5$ m 갔다};
\draw[pathc, line width=0.4pt] (-89, 123) -- (-77, 112);

% ===================== 지금의 배 =====================
% 직사각형이다. 원이 아니다 — 4.9 m x 2.46 m 로 길고 좁다
\draw[plant, line width=1.0pt]
     (-1.23*\S, -2.4*\S) rectangle (1.23*\S, 2.5*\S);
\fill[ink] (0,0) circle (1.4);
% 선수 화살표는 선체 **바깥** 오른쪽에 둔다. 호 21개가 전부 원점에서 뻗어 나오므로
% 가운데에 그리면 호 다발에 묻혀 읽히지 않는다 (2026-10-27 에 한 번 묻혔다)
\draw[axis, -{Stealth[length=5pt,width=4pt]}]
     (1.23*\S+8, 1.4*\S) -- (1.23*\S+8, 2.4*\S);
\node[note, anchor=west] at (1.23*\S+12, 1.9*\S) {$x_b$ 선수};
\node[note, plant, anchor=north, text width=50mm, align=center] at (0,-2.4*\S-4)
     {선체 외곽 $4.9 \times 2.46$ m \\ 원이 아니라 직사각형으로 잰다};

% ===================== 장애물 둘 =====================
\foreach \px/\py in {0.7/6.1, 1.1/5.2, 1.6/5.9, 1.9/5.0, 1.2/6.5, 2.3/5.6}
  {\fill[errc] ({\px*\S},{\py*\S}) circle (1.3);}
\node[note, errc, anchor=west] at (2.3*\S+6, 5.35*\S) {장애물 묶음 1};

\foreach \px/\py in {3.6/2.4, 4.0/3.0, 3.3/3.3, 4.3/2.1}
  {\fill[errc] ({\px*\S},{\py*\S}) circle (1.3);}
\node[note, errc, anchor=west] at (4.3*\S+6, 2.7*\S) {장애물 묶음 2};

% 직진을 지우는 점 하나 — 선체 폭(+-1.23 m) 안에 들어와 있다
\draw[warn, line width=0.6pt] (0.7*\S, 6.1*\S) circle (3.2);
\draw[warn, line width=0.4pt] (0.7*\S+2.6, 6.1*\S+2.0) -- (96, 136);
\node[note, warn, anchor=north west, text width=52mm] at (97, 139)
     {이 점 하나가 선체 폭 $\pm1.23$ m 안에 들어와 있다 \\
      — 직진이 지워지는 이유이고 \texttt{oa\_flag} 가 1 이 되는 이유다};

% ===================== 목표 방향 =====================
\draw[fb] (0,0) -- (2.6*\S, 7.7*\S);
\node[note, guid, anchor=south east] at (2.55*\S, 7.75*\S) {목표점 쪽};

% ===================== 범례 — 그림 아래 =====================
\node[note, warn, anchor=north west, text width=150mm] at (-128, -66)
     {붉은 점선 — 3초 안에 선체가 닿는 호. 점수를 매기지 않고 바로 지운다};
\node[note, anchor=north west, text width=150mm] at (-128, -75)
     {회색 실선 — 살아남은 호. 이 중에서 점수로 고른다};
\node[note, pathc, anchor=north west, text width=150mm] at (-128, -84)
     {굵은 청록 — 고른 호. 선회율 $+12^\circ$/s, 반지름 $7.16$ m};

% ===================== 점수 =====================
\node[anchor=north west] at (-128, -98) {\small \textbf{점수} — 두 가지를 더한다};
\node[anchor=north west] at (-128, -110)
     {$\displaystyle \text{score} \;=\;
       \underbrace{\bigl(1 - \text{align}\bigr)}_{\text{목표를 향하는가}}
       \;+\; 0.3 \cdot
       \underbrace{\text{clip}\!\left(\frac{d_{min}-r}{3},\,0,\,1\right)}_{\text{여유가 넉넉한가}}$};

\node[note, anchor=north west, text width=190mm] at (-128, -134)
     {$\text{align}$ 은 2초 뒤 헤딩이 목표 방위에서 얼마나 벗어나는가를 $0{\sim}1$ 로 잰 것이고,
      $d_{min}$ 은 \textbf{직사각형 선체에서 장애물 점까지의 거리}, $r=0.60$ m 는 안전반지름,
      $3$ m 는 "이만큼 여유가 있으면 만점" 의 기준이다. 가중치가 $0.3$ 이므로
      \textbf{목표 쪽으로 가는 것이 여유보다 세 배 무겁다} — 통로가 열려 있으면 굳이 넓은 쪽으로 돌지 않는다};

% ===================== 주의 =====================
\node[wnote, anchor=north west, text width=190mm] at (-128, -154)
     {\textbf{3초는 $4.5$ m 다 — 자기 배 길이 $4.9$ m 보다 짧다.}
      호 다발 전체가 폭 $2.7$ m $\times$ 깊이 $4.5$ m 밖에 안 되는 손바닥만 한 부채다.
      그래서 이 회피기는 \textbf{멀리 보는 기계가 아니라 가까이를 놓치지 않는 기계}다 —
      라이다를 $130$ m 까지 늘려도 쓰이지 않는다 ($\texttt{lidar.max\_range\_m}$ 기본값이 $10$ m 인 이유).
      실제 제약은 선회율이 $12^\circ$/s 뿐이라는 것이다. 3초에 $36^\circ$ 밖에 못 돈다};

\end{tikzpicture}
\end{document}
